\documentclass[openany,fontset=fandol]{ctexbook}
\usepackage{amsmath,amssymb}
\usepackage[unicode,hidelinks]{hyperref}
\setlength{\emergencystretch}{2em}
\newcommand{\OO}{\mathcal O}
\DeclareMathOperator{\Sym}{Sym}
\DeclareMathOperator{\rk}{rk}
\begin{document}
\title{代数几何基础第五卷预备笔记\\17.16　光滑二次型：16.1--16.3}
\author{据英译编辑稿续译，参校法文扫描}
\date{2026年11月21日}
\maketitle
\chapter*{范围与接续}
本批独立译出§17.16标题及§§16.1--16.3全文：英译物理第78--79页，法文扫描物理第161--162页。止于引理16.3；下一起点在英译物理第79页、法文物理第162页，§16.4：“以下研究主要是为了给出16.2未覆盖情形中的二次型光滑性判据”。不改合订稿。
\chapter*{17.16　光滑二次型}
\noindent\textbf{16.1.}\quad 设$Q$是$\Sym^2(E)$的一个截面，其中$E$为有限秩局部自由模。$Q$也可视为$E^\vee$上的二次型，它定义$E^\vee$上的对称双线性型$B$，从而给出同态$E^\vee\to E$。取行列式模，便得到$\det(E^\vee)\to\det(E)$，最终得到一个截面
\[
d(Q)\in\Gamma\!\left(S,(\det E)^{\otimes 2}\right),
\]
称为二次型$Q$的\emph{判别式}（如果没有弄错的话）。若$E=\OO_S^n$，它就是$\OO_S$的一个截面，即表示这一配对的$n\times n$矩阵的行列式，矩阵系数在$\Gamma(S,\OO_S)$中。先立即指出：

\noindent\textbf{命题 16.2.}\quad 设$E$的秩处处为正，并且对每个$s\in S$，有$\operatorname{char}\kappa(s)\ne2$，或者$E\otimes\kappa(s)$的秩为偶数。则$Q$光滑的充要条件是$Q$“非退化”，即$d(Q)$可逆。事实上，在正秩约定下，非退化总是充分条件，不必对$S$或$E$再加限制。\footnote{原稿未写“秩处处为正”，且充分性原作“不限制$S$或$E$”。正文明确采用正秩。若允许秩零，空矩阵行列式为$1$，而$\mathbf P(0)=V(Q)=\varnothing$，违反15.1中零概形不能等于整条射影纤维的要求。}

最后一句只是便于以后引用：它已包含在前一断言里，因为二次型非退化就蕴含$E$在剩余域特征为$2$的每一点都是偶秩。也可以不区分特征而直接证明：在代数闭域上取一组基，非退化就是说各$\partial Q/\partial X_i$没有共同的非零零点，自然也就没有与$Q$共同的非零零点。若特征不为$2$，反过来也成立；因为
\[
2Q=\sum_i X_i\frac{\partial Q}{\partial X_i},
\]
各偏导数的任何公共零点也都是$Q$的零点。

最后，若特征为$2$，则与$Q$关联的双线性型$B$（其矩阵是各线性型$\partial Q/\partial X_i$的系数矩阵）是交替型，其“核”$N\subseteq E^\vee$满足$E^\vee/N$为偶秩。若$E$也为偶秩，$N$便同样为偶秩。因此，若$N\ne0$，其秩至少是$2$，故在$N$中至少有一个$Q$的非零零点（仍在代数闭域上），即$Q$不光滑。

注意：这样借助坐标实在有失我们的体面；应当只在简单证明下列引理时使用坐标：

\noindent\textbf{引理 16.3.}\quad 设$E$是域$k$上的有限秩向量空间，$Q\in\Sym^2(E)$，$x$是$E^\vee$中的非零向量，$x'$是它在$\mathbf P(E)$中的像。设$Q(x)=0$，即$x'\in V(Q)$。则$V(Q)$在$x'$处光滑且局部维数为$\rk E-2$的充要条件是$x$不属于关联对称双线性型$B$所定义的同态$E^\vee\to E$的核。\footnote{法文及英译都写$Q(x,x)=0$和$\rk E-1$。前者按二次型本身写作$Q(x)=0$，不能换成$B(x,x)=0$（特征$2$时后者对所有$x$成立）。后者改为射影超曲面的维数$\rk E-2$；$\rk E-1$是非零方程所定义的仿射锥的维数。两底本还注明$Q\ne0$；它是维数要求或非零极化微分所蕴含的条件。非退化秩一型的射影零概形为空，故没有引理所设的$x'$。}

\section*{译校说明（非原文）}
\begingroup\small
\paragraph{极化与判别式。}
此处的约定是
\[
B(v,w)=Q(v+w)-Q(v)-Q(w),
\]
没有因子$1/2$；从而在任意剩余特征中均有
$dQ_x(w)=B(x,w)$。矩阵为该配对在一组坐标中的系数矩阵；如果写成$Q=\sum_i a_i X_i^2+\sum_{i<j} b_{ij}X_iX_j$，则其对角项为$2a_i$，非对角项为$b_{ij}$（按对称方式排列）。这里判别式取极化矩阵的行列式。原稿的保留语“如果没有弄错的话”附于名称；英译将它印成脚注72，英译脚注73则来自“立即”一语，二者不是另立的数学命题。

\paragraph{射影切空间如何检验。}
记$W=E^\vee$。在$x'$选取仿射坐标片，用线性函数$\ell$令$\ell(x)=1$（不改变所代表的射影点）。片上$\ell=1$给出线性补空间$\ker\ell$作为切向量空间。因为
\[
B(x,x)=2Q(x)=0,
\]
$B(x,-)$在$kx$上为零，故限制到$\ker\ell$与取商$W/kx$得到同一非零性判据。非零的$Q$在此给出一个超曲面；它在此点光滑且局部维数为$\rk E-2$，恰当且仅当该微分不为零。若$Q=0$，则零概形是整个射影空间，虽光滑但其维数为$\rk E-1$。特征$2$、偶秩且$N\ne0$时，在代数闭域上$Q|_N=\sum a_i X_i^2=(\sum\sqrt{a_i}X_i)^2$；根基维数至少为$2$，这个线性型必有一个非零零点。不能由此断言原域上一定有这样的非零零点。

\paragraph{相对基底。}
按15.1的用语，正秩时先要求$Q$在每一纤维上非零，再要求射影零概形光滑；秩一非零型的空零概形也光滑。每个纤维中非零方程所定义的射影零概形是有效 Cartier 除子（允许为空），于是相对零概形局部有限表示且平坦，见 \href{https://stacks.math.columbia.edu/tag/062Y}{相对有效 Cartier 除子判据}。在这个前提下，再由\href{https://stacks.math.columbia.edu/tag/01V8}{平坦有限表示态射的逐纤维光滑判据}从几何纤维的论证得出相对结论。非退化时，正秩保证各纤维的$Q$不为零。
\par\endgroup
\end{document}
